% !TeX program = xelatex
\documentclass[11pt]{article}

\usepackage[a4paper,margin=2.4cm]{geometry}
\usepackage{fontspec}
\usepackage{microtype}
\usepackage{hyperref}
\usepackage[spanish,es-nodecimaldot]{babel}

\setmainfont{Latin Modern Roman}
\setlength{\parindent}{0pt}
\setlength{\parskip}{0.65em}
\hypersetup{
  colorlinks=true,
  linkcolor=black,
  urlcolor=blue
}

\begin{document}

\begin{center}
{\Large \textbf{La introducción en un plan de tesis}}\\[0.3em]
{\large Función, alcance y forma de redacción}
\end{center}

Luis Koc 

La introducción de un plan de tesis es la puerta de entrada al documento. No está para adornar el tema ni para repetir, capítulo por capítulo, lo que vendrá después. Su tarea es más concreta: ayudar al lector a entender qué situación origina el estudio, por qué merece atención académica o técnica y hacia dónde se dirige la investigación propuesta. Para lograrlo, debe explicar el contexto, identificar una dificultad o brecha, justificar el interés del trabajo y conducir hacia la pregunta de investigación, de acuerdo con la estructura recomendada por Lövei (2021).

El lector previsto suele ser un profesor universitario, un asesor, un jurado, un revisor metodológico o un investigador que necesita evaluar si el plan tiene sentido. Ese lector no busca solo datos; busca una razón clara para seguir leyendo. Quiere ver si el tema es pertinente, si el problema está bien delimitado, si la literatura se usa con criterio y si la propuesta parece viable. Por eso, la introducción debe cuidar el primer contacto con el lector. Después del título y del resumen, esta sección es una de las primeras oportunidades para mostrar la dirección, al claridad y el valor de la investigación.

El inicio debe ser amplio, pero no excevamente genérico o vago. Es conveniente empezar explicando el fenómeno, la situación, el campo de aplicación o la necesidad real que sostiene el estudio. Lövei recomienda una estructura de embudo: iniciar desde un contexto general y avanzar gradualmente hacia el punto específico que será investigado (Lövei, 2021). Una buena apertura no se limita a decir que un tema ``es importante''; muestra en qué contexto aparece, quién se ve afectado, qué cambio reciente lo vuelve relevante o qué dificultad práctica o teórica lo hace digno de estudio.

Después de ese inicio, el texto debe ir cerrando el foco. Primero debe presentar el contexto necesario; luego mostrar la dificultad; finalmente señalar la brecha que justifica la investigación. Scholz sostiene que la introducción debe presentar la motivación, el objetivo del trabajo y la literatura relevante sin repetir conocimiento básico que el público especializado ya domina (Scholz, 2022). En un plan de tesis, esto significa escoger lo indispensable: no se desarrolla todo el marco teórico, pero sí se ofrecen las piezas necesarias para que el lector entienda por qué el estudio debe realizarse.

La problemática puede aparecer en la introducción, pero solo como una primera aproximación. Se presenta como la situación que da origen al estudio, no como una sección completa metida dentro de otra. El problema de investigación también puede anunciarse, aunque su formulación completa corresponde al apartado de planteamiento del problema. Del mismo modo, los objetivos, las variables o categorías principales y las fuentes de evidencia pueden mencionarse si ayudan a entender la línea argumental, pero no deben ocupar el lugar de sus secciones propias.

La metodología debe aparecer como una orientación, no con todo el desarrollo técnico. La introducción puede indicar, de manera general, cómo se abordará el problema: si el estudio analizará casos, construirá un modelo, comparará alternativas, revisará evidencia, aplicará instrumentos o evaluará un artefacto. Lo que no conviene es explicar todavía cada fase, técnica, muestra, métrica o procedimiento. Scholz advierte que la introducción debe conducir al objetivo y al enfoque del estudio, pero sin desplazar a las secciones metodológicas (Scholz, 2022).

Las fuentes de datos y de evidencia siguen la misma lógica. Pueden mencionarse si ayudan a mostrar que el plan es viable y que no se sostiene solo en una impresión personal. Pero la introducción no debería convertirse en una lista de documentos, tablas, archivos o procedimientos. En esta parte basta con que el lector perciba que existe una base de evidencia suficiente; el detalle puede desarrollarse luego en la metodología, el marco teórico o los anexos, según corresponda.

El cierre también es importante. Una buena introducción no termina de golpe ni deja al lector suspendido. Debe preparar el tránsito hacia la problemática, donde se amplía la evidencia; hacia el planteamiento del problema, donde se formula la dificultad central; hacia los objetivos, donde esa dificultad se convierte en dirección de trabajo; hacia el marco teórico, donde se explican los conceptos necesarios; y hacia la metodología, donde se define cómo se evaluará o desarrollará la propuesta. Así, la introducción funciona como puente entre el interés inicial del lector y la arquitectura completa del plan.

El tiempo verbal debe responder a la función de cada idea. El presente sirve para hechos aceptados, definiciones, relaciones vigentes y afirmaciones generales. El pasado se usa para investigaciones, antecedentes o hallazgos ya publicados. El futuro, o una forma impersonal proyectiva, se usa para acciones que el plan todavía propone realizar. Esta combinación es válida si cada tiempo cumple una función clara. Van Emden y Becker (2017) recomiendan una escritura técnica directa, precisa y orientada al lector; esa precisión también exige no presentar como resultado aquello que aún es una intención del estudio.

El nivel de profundidad debe ser selectivo. La introducción no es el marco teórico, porque no desarrolla toda la teoría. Tampoco es la problemática, porque no acumula toda la evidencia. No es la metodología, porque no describe cada técnica en detalle. Y tampoco es el resumen, porque no condensa todo el documento. Su profundidad debe permitir entender qué se estudia, por qué importa, qué brecha existe y cómo se orienta el plan. Lövei (2021) insiste en que la introducción debe preparar al lector para la información nueva, no sustituir el desarrollo completo del estudio.

La extensión no tiene una medida universal. En un plan breve, una introducción puede ocupar una o dos páginas. En un plan con mayor densidad conceptual puede necesitar más espacio. El criterio principal no debe ser el número de hojas, sino la función que cumple el texto. Mientras permita comprender el contexto, la brecha, la relevancia y la dirección general del estudio, la extensión estará bien. Si repite teoría, adelanta resultados, enumera procedimientos o responde preguntas que pertenecen a otras secciones, probablemente ya se extendió más de lo necesario.

Las referencias sí pueden citarse en la introducción cuando sostienen antecedentes, definiciones, datos, brechas o decisiones de enfoque. La cita debe ir cerca de la afirmación que respalda. Puede integrarse en la oración, cuando el autor forma parte del argumento, o ubicarse al final de la frase que contiene la idea sustentada. No conviene colocar una cita al final de un párrafo con varias afirmaciones distintas, porque se pierde trazabilidad. Scholz (2022) recomienda usar literatura relevante para situar el trabajo; esa relevancia exige selección de las fuentes citadas, no acumulación.

Una introducción bien escrita tiene una forma reconocible. Inicia con un contexto significativo, avanza hacia una dificultad concreta, muestra una brecha, justifica el interés del estudio, anuncia la orientación de la propuesta y cierra conectando con el resto del plan. En ese recorrido, la introducción no solo informa. También ordena la lectura y le muestra al lector que el plan tiene una línea argumental clara y defendible.

\vspace{0.5em}
\textbf{Referencias}

Lövei, G. L. (2021). \textit{Writing and Publishing Scientific Papers: A Primer for the Non-English Speaker}. Open Book Publishers. \url{https://doi.org/10.11647/obp.0235}

Scholz, F. (2022). Writing and publishing a scientific paper. \textit{ChemTexts, 8}, Article 8. \url{https://doi.org/10.1007/s40828-022-00160-7}

Van Emden, J., y Becker, L. (2017). \textit{Writing for Engineers} (4th ed.). 
Bloomsbury Academic.
ISBN 9781352000481.

\end{document}
